\chapter{Résilience}\label{ch:ll:resilience}

Le 1\up{er} octobre 2020, un module de mémoire est tombé en panne dans l'un des deux appareils d'un système à disque partagé installé au sein de la Bourse de
Tokyo, dans son système de négociation. Les opérations devaient «~basculer automatiquement~» vers le second appareil, mais la \omterm{def:ll:resilience:failover}{bascule} automatique
ne fonctionnait pas pour ce type de panne, sauf si un paramètre avait été modifié, et il ne l'avait pas été. Les ingénieurs ont basculé
à la main à 9~h~26, moins d'une demi-heure après l'heure d'ouverture prévue du marché, et chaque fonction est redevenue normale. La
négociation est néanmoins restée suspendue toute la journée. Les ordres reçus avant que la bourse ne coupe son réseau avaient été appariés, et
leurs exécutions s'étaient accumulées dans le système sans être transmises aux participants~; seul un nombre limité de
participants aurait pu renvoyer leurs ordres, et il n'existait aucune règle convenue et testée pour le redémarrage. Le matériel a été rétabli en
quelques minutes~; l'état, lui, ne l'a pas été. Ce chapitre traite du second problème~: comment un système de négociation conserve une histoire unique et convenue de ce
qu'il a reçu et envoyé, afin que, lorsqu'une machine tombe en panne, une autre puisse prendre le relais en quelques millisecondes sans perdre un ordre
ni en envoyer un deux fois.

\section{Modes de défaillance d'un système de négociation}

Les composants tombent en panne de quelques façons, et chacune appelle une réponse différente. Un processus plante, ou sa machine plante~: il s'arrête, et la
question est de savoir qui prend le relais et à partir de quel état. Un processus se bloque ou ralentit~: il n'est pas arrêté, il peut donc revenir et
agir après qu'un autre a pris le relais, ce qui est pire qu'un plantage. Une liaison réseau tombe~: chaque côté peut croire que
l'autre est mort alors que les deux sont vivants. Une session avec une plateforme se coupe~: les ordres déjà transmis poursuivent leur course, peuvent être exécutés ou
annulés par celle-ci (chapitre~21), et la société ne sait pas lequel tant qu'elle ne s'est pas reconnectée. Enfin, le mécanisme de \omterm{def:ll:resilience:failover}{bascule} lui-même tombe en panne, comme à
Tokyo, parce qu'il ne s'exécute que lorsque quelque chose d'autre a déjà mal tourné, et c'est donc le code le moins sollicité du
système. Pour un système de négociation, le risque opérationnel du livre~6 (chapitre~28) se résume pour l'essentiel à cette liste.

L'objectif de conception en découle~: chaque entrée sur laquelle le système agit, et chaque ordre qu'il envoie, doit être enregistrée avant de produire un effet, dans
un ordre unique que tous acceptent, de sorte que n'importe quelle machine puisse reconstruire l'état en rejouant l'enregistrement, et que tout ordre puisse
être reconnu s'il est renvoyé.

\section{Architectures à séquenceur}

\begin{definition}[Architecture à séquenceur]\label{def:ll:resilience:sequencer}
Une \emph{architecture à séquenceur}\index{architecture à séquenceur} est une conception dans laquelle chaque entrée d'un système (événements de marché,
comptes rendus d'ordres, commandes) passe par un composant unique, le séquenceur, qui lui attribue le numéro suivant d'une séquence globale
et l'ajoute à un journal durable avant que tout autre composant n'agisse sur elle~; chaque composant traite ensuite le
journal dans l'ordre des numéros.
\end{definition}

Le séquenceur transforme de nombreuses sources non ordonnées en une histoire ordonnée unique. C'est le \omterm{def:ll:logging-capture-and-replay:journal}{journal des entrées} du chapitre~23, élargi du
processus au système~: les données de marché venant du \omterm{def:ll:the-feed-handler:handler}{feed handler}, les comptes rendus venant de la \omterm{def:ll:order-gateway-and-order-management:gateway}{passerelle d'ordres}, les changements de paramètres venant du
plan de contrôle deviennent tous des entrées numérotées d'un même journal. Le rejeu du moteur de stratégie, présenté au chapitre~23, devient le mode
de fonctionnement normal~: un composant est une fonction du journal, et deux composants qui lisent le même journal calculent la
même chose. Le coût est un saut supplémentaire sur le chemin de chaque entrée (une microseconde dans la simulation de la construction, davantage lorsque le
journal est répliqué avant d'être acquitté), ainsi qu'un composant dont la défaillance doit être gérée elle-même~; la bourse de détail dont Fowler a décrit
l'architecture (chapitre~20) gardait sa logique métier sur un seul fil d'exécution alimenté par un journal de ce type, à partir duquel son état pouvait être reconstruit.

\begin{omfigure}
\begin{tikzpicture}[font=\scriptsize,>=Stealth,box/.style={draw,rounded corners=2pt,minimum height=0.7cm,align=center,
  inner sep=3pt,fill=omDef!12}]
\node[box] (md) at (0,1.0) {données de marché\\ (feed handler)};
\node[box] (rp) at (0,-1.0) {comptes rendus\\ (passerelle)};
\node[box,fill=omThm!18] (sq) at (3.0,0) {séquenceur\\ numérote les entrées};
\node[box,fill=gray!10,minimum width=2.2cm] (jr) at (6.0,0) {journal\\ époque $|$ séq $|$ entrée};
\node[box] (pr) at (9.2,1.0) {réplique primaire\\ \textbf{envoie}};
\node[box] (bk) at (9.2,-1.0) {réplique de secours\\ calcule, sans envoyer};
\node[box,fill=gray!10] (vn) at (12.3,1.0) {plateforme};
\draw[->] (md) -- (sq);
\draw[->] (rp) -- (sq);
\draw[->,thick] (sq) -- (jr);
\draw[->] (jr) -- (pr);
\draw[->] (jr) -- (bk);
\draw[->,thick] (pr) -- node[above]{ordres} (vn);
\draw[->,dashed] (pr) -- node[right]{heartbeats} (bk);
\draw[->,dashed,omThm] (bk.east) -- ++(0.5,0) |- node[pos=0.25,right,align=left]{bascule~: nouvelle\\ époque, reconnexion,\\ réconciliation, renvoi} (vn.south);
\end{tikzpicture}

\omcaption{Une \omterm{def:ll:resilience:sequencer}{architecture à séquenceur} avec un secours à chaud. Chaque entrée est numérotée et journalisée avant que quiconque n'agisse sur elle~; les deux répliques appliquent le journal et calculent le même état et les mêmes ordres, mais seule la
réplique primaire envoie. Lorsque les \omterm{def:ll:protocols-i-fix:session}{heartbeats}
cessent, le secours élève l'époque du journal, rétablit la session de saisie des ordres, réconcilie et renvoie ce que
aucun compte rendu n'acquitte.}\label{fig:ll:resilience:arch}
\end{omfigure}

\section{Machines à états répliquées et bascule}

\begin{definition}[Machine à états répliquée, réplication primaire--secours]\label{def:ll:resilience:rsm}
Une \emph{machine à états répliquée}\index{machine à états répliquée} est un programme déterministe exécuté en plusieurs copies (répliques)
qui appliquent la même suite d'entrées et passent donc par les mêmes états et produisent les mêmes sorties.
La \emph{réplication primaire--secours}\index{réplication primaire--secours} fait fonctionner une réplique comme primaire, dont les sorties sont envoyées,
et les autres comme secours, dont les sorties sont calculées puis écartées, prêtes à remplacer le primaire.
\end{definition}

\begin{definition}[Bascule, split brain]\label{def:ll:resilience:failover}
Une \emph{bascule}\index{bascule} (failover) est le remplacement d'un primaire défaillant par un secours, qui reprend son
rôle et ses connexions. Le \emph{split brain}\index{split brain} est l'état dans lequel deux répliques agissent simultanément comme primaire, typiquement
parce qu'un secours a pris le relais d'un primaire qui était lent ou coupé du réseau plutôt que mort~; chacune envoie alors
des sorties dont l'autre ne sait rien.
\end{definition}

L'approche par machine à états est ancienne et générale (le tutoriel de Schneider, en 1990), et sa partie difficile, c'est-à-dire se mettre d'accord sur
l'entrée suivante du journal lorsque des machines tombent en panne, est ce que résolvent des algorithmes de consensus tels que Raft (Ongaro et Ousterhout, 2014). Un système de négociation
a généralement besoin de moins~: un séquenceur qui écrit le journal, répliqué de façon synchrone sur une seconde machine, et une règle pour
décider qui est primaire. Les répliques de la construction sont l'état de stratégie du livre en miniature~: la réplique de \texttt{firm.sequencer}
négocie cent actions à chaque mouvement de prix avec un ordre exécuter-ou-annuler, dans une limite de position, et conserve
ses ordres au moyen d'un identifiant dérivé de l'entrée du journal qui les a créés. Les répliques en Python, en C++ et en Rust rejouent
le journal de l'exemple de référence (une \omterm{def:ll:resilience:failover}{bascule}, avec une seconde époque, des renvois et un doublon) et atteignent les mêmes états, empreinte de hachage comprise.

Le \omterm{def:ll:resilience:failover}{split brain} est empêché par le \emph{fencing}. Le journal porte une époque, un nombre qui ne fait que croître~; chaque primaire écrit sous
l'époque dans laquelle il est devenu primaire, et un secours qui prend le relais en premier relève l'époque du journal. Dès lors, le journal
refuse toute écriture sous une époque antérieure (\cref{lst:ll:resilience:fence})~: un primaire qui n'était que lent, et qui se réveille en se croyant encore
aux commandes, peut calculer ce qu'il veut mais ne peut plus rien enregistrer, et rien de ce qui n'a pas été enregistré
n'est envoyé. La passerelle applique la même règle aux sorties, puisque la plateforme ignore les époques.

\omcode{code/firm/sequencer/cpp/firm_sequencer.hpp}{32}{42}{La barrière du journal~: après une nouvelle époque, une écriture issue d'une époque antérieure
ne peut plus rien ajouter.}\label{lst:ll:resilience:fence}

La défaillance est détectée au moyen des \omterm{def:ll:protocols-i-fix:session}{heartbeats}~: le primaire en envoie une par milliseconde, et le secours déclare mort le primaire après trois
signaux manqués. La détection représente donc la plus grande part du temps de reprise de la construction~: de 2.8 à \qty{3.1}{\milli\second} depuis la panne jusqu'à un secours devenu primaire, reconnecté et réconcilié,
dont la reconnexion et le rejeu des comptes rendus manqués prennent
quelques dizaines de microsecondes. Un délai d'expiration plus court accélère la reprise et déclare mort davantage de primaires lents, ce qui est précisément le cas où le fencing
compte. Le secours de la construction est à chaud~: il applique chaque entrée du journal dès qu'elle est écrite, si bien qu'il n'y a rien à rattraper. Un
secours à froid doit d'abord rejouer le journal~; la réplique C++ l'applique à environ \qty{37}{\nano\second} par entrée, de sorte qu'un journal
d'une journée de 37~millions d'entrées (chapitre~23) prend environ une seconde et demie, acceptable pour un redémarrage, mais pas pour une \omterm{def:ll:resilience:failover}{bascule}.

\section{Coupures de l'échange et reprise}

\begin{definition}[Traitement idempotent]\label{def:ll:resilience:idempotent}
Une opération est traitée \emph{de façon idempotente}\index{traitement idempotent} lorsque l'effectuer deux fois a le même effet que
l'effectuer une seule fois. Pour les ordres, cela signifie que chaque ordre porte un identifiant fixé à sa création, de sorte que la plateforme, ou
la passerelle propre à la société, reconnaît une seconde copie et la rejette.
\end{definition}

Le moment dangereux d'une \omterm{def:ll:resilience:failover}{bascule} tient aux ordres de l'intervalle~: ceux que l'ancien primaire a envoyés et dont les comptes rendus n'ont jamais été journalisés.
Le nouveau primaire sait qu'ils existent, puisqu'il les a calculés à partir du journal, mais ignore s'ils sont partis~: l'ancien primaire peut être mort
avant de les envoyer, alors que l'ordre était sur le fil, ou après que la plateforme l'a exécuté et avant que le compte rendu ne revienne. Il
dispose de trois manières de procéder (\cref{lst:ll:resilience:resend}). Il peut les renvoyer aussitôt sous de nouveaux identifiants, réflexe
des passerelles qui ne réutilisent jamais un identifiant. Il peut d'abord se reconnecter, demander à la plateforme les comptes rendus qu'il a manqués (le rejeu séquencé
du chapitre~21), réconcilier, et ne renvoyer que ce qui reste sans trace. Ou bien il peut renvoyer aussitôt sous les identifiants d'origine des ordres
d'origine, dérivés du journal, et laisser la plateforme rejeter toute copie d'un ordre qu'elle possède déjà.

\omcode{code/firm/sequencer/firm_sequencer_sim.py}{116}{137}{La reprise~: barrière du journal, puis renvoi de ce qu'aucun compte rendu
n'accuse, sous les identifiants d'origine ou sous de nouveaux (un renvoi est journalisé, afin que les répliques
restent en accord).}\label{lst:ll:resilience:resend}

Le \Cref{tab:ll:resilience:cuts} rend compte d'un primaire tué à chaque étape de chacun des 46 ordres du scénario, face au
moteur d'appariement du livre~10. Avec de nouveaux identifiants et sans réconciliation, une panne pendant qu'un ordre est sur le fil ou chez la plateforme, mais pas
encore signalé, l'envoie deux fois~: 54 des 184 points de coupure se terminent par un doublon, un pour chacun des 27 ordres exécutés, dans chacune
des deux étapes, et le pire fait dépasser la position à 1\,100 actions, contre 1\,000. Réconcilier d'abord
les supprime tous dans ce modèle, car chaque ordre est signalé en moins de \qty{40}{\micro\second} et la plateforme rejoue ses
comptes rendus. Les identifiants idempotents les suppriment quel que soit l'ordre des opérations, ce qui compte lorsque la plateforme tarde
à signaler ou que la session ne peut pas être rejouée~: les participants de la bourse de Tokyo en 2020 se trouvaient précisément dans cette situation. Dans
tous les cas, l'état du nouveau primaire est égal à un rejeu frais du journal et à la position que détient la plateforme.

\begin{table}[ht]
\centering\small
\begin{tabular}{@{}>{\raggedright\arraybackslash}p{4.0cm}>{\centering\arraybackslash}p{1.9cm}>{\centering\arraybackslash}p{1.9cm}>{\centering\arraybackslash}p{1.9cm}>{\centering\arraybackslash}p{1.9cm}>{\centering\arraybackslash}p{2.3cm}@{}}
\toprule
procédure & \multicolumn{4}{c}{doublons, selon l'étape de l'ordre au moment de la panne} & position maximale \\
 & journalisé & en vol & chez la plateforme & signalé & \\
\midrule
identifiants frais & 0 & 27 & 27 & 0 & 1\,100 \\
identifiants frais, réconciliés d'abord & 0 & 0 & 0 & 0 & 1\,000 \\
clés d'idempotence & 0 & 0 & 0 & 0 & 1\,000 \\
\bottomrule
\end{tabular}
\caption{Le primaire tué à chacune des quatre étapes de chacun des 46 ordres du scénario (184 \omterm{def:ll:resilience:failover}{bascules} par procédure) face au
moteur d'appariement~; la reprise a pris de 2.8 à \qty{3.1}{\milli\second} dans chaque cas, et l'état était toujours égal à un rejeu du
journal. Données~: \texttt{fig\_failover.py} (déterministe).}\label{tab:ll:resilience:cuts}
\end{table}

\begin{omfigure}
\begin{tikzpicture}[font=\scriptsize,>=Stealth]
\draw[->] (0,0) -- (12.6,0) node[right]{temps};
\foreach \x/\l in {0/{entrée journalisée},2.6/{ordre envoyé},6.2/{sur la plateforme},9.8/{compte rendu journalisé}} {
  \draw (\x,-0.12) -- (\x,0.12);
  \node[below,align=center] at (\x,-0.15) {\l};
}
\fill[omDef!15] (0,0.2) rectangle (2.6,1.0);
\node[align=center,font=\tiny] at (1.3,0.6) {non envoyé~:\\ un renvoi\\ est sans risque};
\fill[omThm!25] (2.6,0.2) rectangle (9.8,1.0);
\node[align=center] at (6.2,0.6) {une panne ici~: un renvoi sous un nouvel identifiant\\ duplique l'ordre s'il a été exécuté};
\fill[omMeth!15] (9.8,0.2) rectangle (12.4,1.0);
\node[align=center] at (11.1,0.6) {connu~:\\ rien à renvoyer};
\draw[<->] (2.6,1.25) -- node[above]{\qty{20}{\micro\second} sur le fil, \qty{20}{\micro\second} au retour} (9.8,1.25);
\end{tikzpicture}

\omcaption{La vie d'un ordre et la fenêtre dans laquelle une \omterm{def:ll:resilience:failover}{bascule} peut le dupliquer~: du moment où il part au moment où
son compte rendu est journalisé, deux traversées du réseau. Avec de nouveaux identifiants et sans réconciliation, une panne
n'importe où dans la fenêtre rouge envoie l'ordre deux fois s'il avait été exécuté.}\label{fig:ll:resilience:window}
\end{omfigure}

Le même raisonnement s'applique à une simple déconnexion, sans aucune \omterm{def:ll:resilience:failover}{bascule}~: la passerelle perd sa session, se reconnecte en demandant
les comptes rendus qu'elle n'a pas vus, les applique (l'annulation à la déconnexion de la plateforme a pu entre-temps annuler des ordres au repos),
réconcilie avec le \omterm{def:ll:order-gateway-and-order-management:dropcopy}{drop copy}, et ce n'est qu'alors qu'elle laisse la stratégie négocier. Le simulateur du livre~10 fournit chaque élément~: rejeu séquencé
à la connexion, annulation à la déconnexion, et rejet d'un identifiant d'ordre client en double.

\section{Tester la bascule}

Le code de \omterm{def:ll:resilience:failover}{bascule} s'exécute lorsque quelque chose d'autre a déjà échoué, si bien que c'est le code le moins testé d'un système, à moins qu'on ne le teste
délibérément, et l'incident de Tokyo est l'exemple type d'un dispositif qui a fonctionné pour les pannes contre lesquelles il avait été testé, et non
pour celle qui s'est produite. Le test de la construction est le tableau~: le primaire est tué à chaque étape de chaque ordre d'un scénario,
et chaque exécution est vérifiée contre des propriétés qui doivent tenir quel que soit le point de panne (l'état du nouveau primaire est égal à un rejeu
du journal~; sa position est celle de la plateforme~; aucun ordre n'est exécuté deux fois lorsque des clés d'idempotence sont utilisées). Les propriétés
comptent davantage que le scénario~; des points de coupure exhaustifs transforment un événement rare en test courant. En production, la même idée est
pratiquée en provoquant délibérément des \omterm{def:ll:resilience:failover}{bascules}, selon un calendrier, afin que la procédure soit exercée lorsque rien d'autre ne va mal.

\section{Tutoriel~: tuer le primaire partout}\label{tut:ll:resilience:tut}

\begin{tutorial}
\textbf{Objectif.} Lancer le séquenceur, deux répliques et une passerelle face au moteur d'appariement, tuer le primaire à chaque étape
de chaque ordre, et comparer trois procédures de reprise. \textbf{État final~:} \cref{tab:ll:resilience:cuts}, et des rejeux verts du
journal de \omterm{def:ll:resilience:failover}{bascule} dans trois langages.
\begin{tutsteps}
\item \textbf{Le scénario.} \texttt{firm\_sequencer\_sim.run(None)} exécute 60 cotations, à une milliseconde d'intervalle, sans panne~: 46
ordres, une position de $-700$, l'état égal à un rejeu.
\item \textbf{Points de coupure.} \texttt{cut\_points()} liste le milieu de chacune des quatre étapes de chaque ordre~;
\texttt{python fig\_failover.py} exécute les trois procédures à chacun de ces points.
\item \textbf{Répliques dans trois langages.} Les répliques C++ et Rust rejouent \texttt{data/journal.txt}, une \omterm{def:ll:resilience:failover}{bascule} avec de nouveaux
identifiants, et atteignent les états de \texttt{data/expected.txt}~; un écrivain de l'ancienne époque est refusé.
\item \textbf{Rattrapage.} \texttt{python bench\_sequencer.py} mesure le temps qu'une réplique met à appliquer deux millions d'entrées.
\end{tutsteps}
\textbf{Pour aller plus loin.} Retarder les comptes rendus de la plateforme d'une milliseconde et observer les doublons réapparaître même avec la réconciliation~;
raccourcir le délai d'expiration des \omterm{def:ll:protocols-i-fix:session}{heartbeats} et compter les \omterm{def:ll:resilience:failover}{bascules} qu'un primaire lent provoquerait.
\end{tutorial}

\section{Construction~: le séquenceur}\label{bld:ll:resilience:sequencer}

\begin{build}
\bfield{Rôle} Le backbone qui rend redémarrable et réplicable le chemin de négociation du livre (chapitres~18 à 22)~: un
journal unique d'entrées numérotées, des répliques qui calculent à partir de lui, et une \omterm{def:ll:resilience:failover}{bascule} qui ne perd ni ne duplique aucun ordre.
\bfield{Interface} Référence Python \texttt{firm\_sequencer}~: \texttt{Journal} avec \texttt{append} et \texttt{fence}~;
\texttt{Replica} avec \texttt{apply} et \texttt{unacknowledged}~; \texttt{cl\_for}~; \texttt{to\_line}, \texttt{from\_line}~;
\texttt{firm\_sequencer\_sim.run} et \texttt{cut\_points} face au moteur du livre~10. C++20 \texttt{firm::seq}~:
\texttt{Journal}, \texttt{Replica}, \texttt{parse}. Rust \texttt{firm\_sequencer}~: \texttt{Journal}, \texttt{Replica},
\texttt{parse}.
\bfield{Règles} Rien n'agit sur une entrée avant qu'elle ne soit journalisée~; les répliques sont des fonctions déterministes du journal~; les époques
ne font que croître et interdisent aux écrivains plus anciens~; les identifiants d'ordre dérivent du journal~; les renvois sont journalisés.
\bfield{Tests d'acceptation} \texttt{code/firm/sequencer/}~: le journal de \omterm{def:ll:resilience:failover}{bascule} rejoué jusqu'aux mêmes états dans trois langages~;
la barrière refuse une époque ancienne~; les trous sont refusés~; à chaque point de coupure du scénario et pour chaque procédure, l'état est égal à un rejeu
et à la position de la plateforme~; les doublons n'apparaissent qu'avec de nouveaux identifiants et sans réconciliation.
\bfield{Pour aller plus loin} Un journal répliqué avec accusé de réception synchrone par le secours~; plusieurs passerelles derrière les répliques~;
Raft pour élire le primaire~; la passerelle du chapitre~21 et la \omterm{def:ll:pre-trade-risk-and-kill-switches:gate}{barrière de risque} du chapitre~22 en tant que répliques.
\end{build}

\begin{omsources}
\item Tokyo Stock Exchange, \emph{The Cause of the Recent Failure in arrowhead and Measures Implemented}, 5 October 2020, et
\emph{The Failure of Equity Trading System on October 1, 2020}, rapport, 19 octobre 2020.
\item F.~B.~Schneider, ``Implementing fault-tolerant services using the state machine approach: a tutorial'', \emph{ACM
Computing Surveys} 22(4), 1990.
\item D.~Ongaro et J.~Ousterhout, ``In search of an understandable consensus algorithm'', \emph{USENIX Annual Technical
Conference}, 2014.
\end{omsources}

\section{Exercices}

\begin{exercise}[$\star$]\label{exo:ll:resilience:1}
L'époque du journal est 3. L'ancien primaire, devenu primaire à l'époque 2, tente d'ajouter une entrée. Que se passe-t-il, et que devient
les ordres qu'il a calculés à partir de cette entrée~?
\end{exercise}

\begin{exercise}[$\star$]\label{exo:ll:resilience:2}
Des battements toutes les millisecondes, trois battements manqués déclarent le primaire mort, vérifiés à chaque battement. Quels sont les temps
les plus courts et les plus longs entre une panne et sa détection~?
\end{exercise}

\begin{exercise}[$\star$]\label{exo:ll:resilience:3}
Un ordre créé par l'entrée 46 du journal est la première sortie de cette entrée. Quel est son identifiant client dans la construction, et
pourquoi chaque réplique doit-elle le calculer de la même façon~?
\end{exercise}

\begin{exercise}[$\star\star$]\label{exo:ll:resilience:4}
Une stratégie envoie 500 ordres par seconde~; chacun passe \qty{40}{\micro\second} entre son départ et l'journalisation de son compte rendu,
et 90~\% d'entre eux s'exécutent. Une panne survient à un instant aléatoire. Quelle est la probabilité qu'un renvoi avec de nouveaux identifiants,
sans réconciliation, duplique un ordre~?
\end{exercise}

\begin{exercise}[$\star\star$]\label{exo:ll:resilience:5}
Une sauvegarde froide applique le journal à raison de \qty{37}{\nano\second} par entrée. Combien de temps lui faut-il pour rattraper une journée de
37 millions d'entrées, puis la dernière heure de celle-ci~? Pourquoi garder malgré tout la sauvegarde à chaud~?
\end{exercise}

\begin{exercise}[$\star\star$]\label{exo:ll:resilience:6}
La plateforme rejette une copie d'un ordre sous un identifiant qu'elle a déjà vu. La sauvegarde renvoie sous les identifiants d'origine un ordre
que la plateforme a exécuté. Que reçoit la sauvegarde, et que doit faire sa réplique de cette réponse~?
\end{exercise}

\begin{exercise}[$\star\star\star$]\label{exo:ll:resilience:7}
\emph{Programmation.} Retarder les comptes rendus de la plateforme d'une milliseconde dans \texttt{firm\_sequencer\_sim} et relancer les trois procédures.
Laquelle duplique maintenant des ordres, à quelles étapes, et pourquoi les clés d'idempotence n'en produisent-elles toujours aucun~?
\end{exercise}

\begin{exercise}[$\star\star\star$]\label{exo:ll:resilience:8}
\emph{Trouver la faille.} «~Notre \omterm{def:ll:resilience:failover}{bascule} est simple~: la sauvegarde surveille le battement du primaire et, quand il s'arrête, elle
relance la stratégie à partir de sa configuration et envoie de nouvelles cotations.~»
\end{exercise}

\section{Problème~: La bascule qui a doublé les ordres}

\begin{problem}[{Problème du week-end --- compter les doublons en chaque point de défaillance}]\label{pb:ll:resilience:1}
Utiliser le \cref{tab:ll:resilience:cuts} (\texttt{failover.csv}, \texttt{failover\_summary.csv})~: 46 ordres en 60 millisecondes, quatre
étapes chacun, trois procédures de reprise.

\medskip\noindent\textbf{Partie I --- La fenêtre.}
\begin{enumerate}
\item Nommer les quatre étapes de la vie d'un ordre utilisées comme points de coupure, et dire pour lesquelles le nouveau primaire ne peut pas savoir
si l'ordre est parti.
\item Combien de temps un ordre passe-t-il dans les deux étapes dangereuses de la simulation~?
\item Pourquoi seuls 27 des 46 ordres produisent-ils un doublon dans chaque étape dangereuse~?
\item Pourquoi une panne survenant alors qu'un ordre est journalisé mais pas encore envoyé ne le duplique-t-elle jamais~?
\end{enumerate}

\medskip\noindent\textbf{Partie II --- Dans le temps.}
\begin{enumerate}[resume]
\item Quel est le taux d'ordres du scénario~?
\item Si l'instant de la panne est uniforme sur le scénario, quelle est la probabilité qu'une \omterm{def:ll:resilience:failover}{bascule} avec de nouveaux identifiants duplique
un ordre~?
\item Quelle serait-elle pour une stratégie envoyant 5\,000 ordres par seconde avec les mêmes latences~?
\item Pourquoi la plus grande position atteint-elle 1\,100 avec de nouveaux identifiants~?
\end{enumerate}

\medskip\noindent\textbf{Partie III --- Les remèdes.}
\begin{enumerate}[resume]
\item Pourquoi réconcilier d'abord supprime-t-il tous les doublons dans ce modèle~?
\item Donner deux situations réelles dans lesquelles réconcilier d'abord ne suffirait pas.
\item Pourquoi les clés d'idempotence fonctionnent-elles quel que soit l'ordre des opérations~?
\item Que doit-il être vrai pour que les clés d'idempotence soient possibles~?
\end{enumerate}

\medskip\noindent\textbf{Partie IV --- Le verdict.}
\begin{enumerate}[resume]
\item Énoncer le \emph{résultat nommé}~: les doublons après une \omterm{def:ll:resilience:failover}{bascule} sans clés d'idempotence sur les points de panne, avec elles,
et le temps de reprise.
\item Qu'est-ce qui domine le temps de reprise, et quel est le prix de le raccourcir~?
\item Qu'apporte le verrouillage par époque que les clés d'idempotence n'apportent pas~?
\item Qu'est-ce qui a mal tourné à Tokyo en 2020, dans les termes du chapitre~?
\item Pourquoi les renvois eux-mêmes doivent-ils être journalisés~?
\item Quelles propriétés le test vérifie-t-il en chaque point de coupure, et pourquoi des propriétés plutôt que des sorties attendues~?
\item Qu'ajouteriez-vous pour tester une sauvegarde froide~?
\item En une phrase~: qu'est-ce qui rend une \omterm{def:ll:resilience:failover}{bascule} sûre~?
\end{enumerate}
\end{problem}

\section{Questions d'entretien}

\begin{interviewq}[$\star$ \iqroles{développeur}]\label{iq:ll:resilience:1}
Qu'est-ce qu'un séquenceur, et pourquoi un système de trading en placerait-il un devant tout le reste~?
\end{interviewq}

\begin{interviewq}[$\star\star$ \iqroles{développeur}]\label{iq:ll:resilience:2}
Votre primaire cesse d'envoyer ses battements. Comment la sauvegarde prend-elle le relais sans envoyer d'ordres en double~?
\end{interviewq}

\begin{interviewq}[$\star\star$ \iqroles{développeur}]\label{iq:ll:resilience:3}
Qu'est-ce que le \omterm{def:ll:resilience:failover}{split brain}, et comment l'empêchez-vous~?
\end{interviewq}

\begin{interviewq}[$\star\star$ \iqroles{développeur, trader}]\label{iq:ll:resilience:4}
Votre session vers la bourse tombe pendant cinq secondes. Que faites-vous, dans l'ordre, quand elle revient~?
\end{interviewq}

\begin{interviewq}[$\star\star$ \iqroles{développeur}]\label{iq:ll:resilience:5}
Comment testeriez-vous que la \omterm{def:ll:resilience:failover}{bascule} fonctionne, sachant qu'elle est rarement nécessaire~?
\end{interviewq}

\begin{interviewq}[$\star\star\star$ \iqroles{développeur}]\label{iq:ll:resilience:6}
Concevez un système de trading qui survit à la perte de n'importe quelle machine, avec au plus quelques millisecondes d'interruption et aucun
ordre perdu ni dupliqué. Quels en sont les composants, le protocole entre eux, et les tests~?
\end{interviewq}
